Biological neural systems must be fast but are energy-constrained. Evolution's solution: act on the first signal. Winner-take-all circuits and time-to-first-spike coding implicitly treat \emph{when} a neuron fires as an expression of confidence.

% Sources:
%   97.2% halt-first: x02.log, x03.log, x04.log (3 seeds × 4 rotations)
%   86.1% baseline: x00.log (unfuck.txt)
%   97.3% TTA: x00.log with 8 rotations (unfuck.txt)
We apply this principle to ensembles of Tiny Recursive Models (TRM)~\citep{jolicoeur2025trm}. By basing the ensemble prediction solely on the first to halt rather than averaging predictions, we achieve 97.2\% puzzle accuracy on Sudoku-Extreme while using 10$\times$ less compute than test-time augmentation (the baseline achieves 86.1\% single-pass, 97.3\% with TTA). Inference speed is an implicit indication of confidence.

% Source: x182{a,b,c,d}.log (randn): 96.85% ± 0.58%
But can this capability be manifested as a training-only cost? Evidently yes: by maintaining $K{=}4$ parallel latent states during training but backpropping only through the lowest-loss ``winner,'' a single model achieves \textbf{96.9\% $\pm$ 0.6\%} puzzle accuracy---roughly matching TTA performance without any test-time augmentation.

As in nature, this work was also resource constrained: all experimentation used a single RTX 5090. This necessitated efficiency and compelled our invention of a modified SwiGLU~\citep{shazeer2020swiglu} which made Muon~\citep{jordan2024muon} viable. With these improvements and $K{=}1$ training, we match TRM baseline performance ($\sim$87\%) in just 48 min (8k steps, bs${=}$384, ${\sim}$16GiB). Higher accuracy ($\sim$97\%) is achieved in 36k steps and $K{=}4$ (bs${=}$192, ${\sim}$30GiB) and takes about 6 hours.

% The path: biological principle $\to$ inference strategy that validates it $\to$ training procedure that internalizes the benefit. Speed is confidence.
